\question

Dada la integral
\begin{math}\displaystyle
    I=
    \int_{0}^{1}
    \int_{0}^{x}
    \left(1-x^{2}\right)
    \dl y\dl x
\end{math}.

\begin{parts}
    \part[5]

    Describa geométricamente el sólido que representa dicha
    integral doble.

    \part[5]

    Exprese dicha integral realizando un cambio de integración, y
    luego determine el valor de dicha integral.
\end{parts}

\begin{figure}[ht!]
    \centering
    \begin{tikzpicture}[xscale=0.5,yscale=0.5]
        \draw[help lines] (0,0) grid (32,12);
    \end{tikzpicture}
\end{figure}

\begin{solutionordottedlines}
    \begin{parts}
        \part

        \textbf{Descripción geométrica}.


        \begin{multicols}{2}
            La integral doble representa el volumen bajo la
            superficie $f\left(x,y\right)=1-x^{2}$ sobre la
            región triangular en el plano $xy$, donde
            $0\leq y\leq x\leq 1$.
            Es un sólido con base triangular y altura decreciente
            cuadráticamente en $x$.
            La línea roja $y=x$ marca el límite superior de $y$
            en la región de integración.
            La línea azul $f\left(x,y\right)=1-x^{2}$ muestra
            cómo decrece la altura de la superficie en el plano
            $x$ versus $f\left(x,y\right)$.

            \begin{center}
                \includegraphics[width=.3\paperwidth]{2}
            \end{center}

        \end{multicols}

        \part

        \textbf{Cambio de orden de integración y cálculo}.

        La función $f\left(x,y\right)=1-x^{2}$ es continua
        en $\left\{\left(x,y\right)\mid 0\leq y\leq x\leq 1\right\}$
        que es un conjunto compacto, así es $f\left(x,y\right)$
        es integrable.
        La integral doble converge absolutamente, ya que
        $\left|1-x^{2}\right|\leq 1$ en $\left[0,1\right]$.
        \begin{align*}
            \shortintertext{Integral original:}
            I & =
            \int_{0}^{1}
            \int_{0}^{x}
            \left(1-x^{2}\right)
            \dl y\dl x.
            \shortintertext{Cambiando el orden $0\leq y\leq 1$, $y\leq x\leq 1$ (Teorema de Fubini):}
            I & =
            \int_{0}^{1}
            \int_{y}^{1}
            \left(1-x^{2}\right)
            \dl x\dl y.
            \shortintertext{Integrando respecto a $x$:}
            I & =
            \int_{0}^{1}
            \left[
            \left.
            \left(x-\dfrac{x^{3}}{3}\right)
            \right|^{x=1}_{x=y}
            \right]\dl y=
            \int_{0}^{1}
            \left[
                \left(1-\dfrac{1^{3}}{3}\right)-
                \left(y-\dfrac{y^{3}}{3}\right)
                \right]\dl y=
            \int_{0}^{1}
            \left(
            \frac{2}{3}-y+\frac{y^{3}}{3}
            \right)\dl y.
            \shortintertext{Integrando respecto a $y$:}
            I & =
            \left.
            \left(
            \dfrac{2}{3}y-\dfrac{y^{2}}{2}+\dfrac{y^{4}}{12}
            \right)
            \right|_{y=0}^{y=1}=
            \dfrac{3}{12}.
        \end{align*}
    \end{parts}
\end{solutionordottedlines}
